\section{Discussion and Scope}
\label{sec:limitations}

UWM-JEPA is a JEPA-family world model whose latent is a density matrix
and whose predictor is unitary on a system--environment space. This
construction yields an exact non-dissipation property on the joint
latent and, in the tested partially observable settings, improves
intermediate-horizon blind rollout and counterfactual
action-conditioned imagination. The architecture is not a universally
stronger encoder, a general anti-collapse mechanism, or a competitive
planner.

The phrase \emph{belief state} is used here in an operational sense:
the latent is a density-matrix representation capable of carrying
uncertainty and system--environment correlations through rollout.
Calibrated posterior structure over the true hidden state would
require a separate posterior-alignment experiment and is outside the
present claims.

The headline negative result is the context-probe tie. UWM-JEPA and
LSTM-JEPA encode the observed prefix equally well on the held-out
hidden-velocity probe, so the empirical advantage appears only when
the latent is rolled forward without new observations or when actions
are varied counterfactually. This separation organises the experiments
around encoder quality, blind rollout, and action binding rather than
around a single aggregate score.

The non-forgetting theorem applies to the joint density matrix.
Unitary conjugation preserves the spectrum, purity, and entropy of the
\emph{joint} state, but the downstream loss and probes see only the
reduced and projected system state. Reduced-state information can
move into correlations and become inaccessible to a linear probe even
when the joint evolution is exactly reversible, so the operational
memory curves in \cref{fig:blind-rollout} complement rather than
duplicate the theorem; they measure what remains useful after the
readout used by the actual JEPA objective. The JEPA predictors in
that experiment are trained at horizon $k=5$, so $k=1$ and $k=3$
probe intermediate rollout coherence and $k>5$ probes extrapolation
rather than a monotone information-theoretic forgetting curve.

The action results establish counterfactual sensitivity rather than
broad planning competence. Teacher-forced action JEPA leaves the action
Hamiltonian nearly inert, whereas counterfactual targets recover an
action term comparable to the base dynamics and yield an above-chance
hidden-velocity indicator with the expected degradation under wrong,
shuffled, and no-action controls. The construction is evidence for
action-conditioned imagination in this testbed, not for long-horizon
planning, and not for alternative recurrent or Transformer objectives.
The counterfactual training target also requires simulator access to
construct alternate futures, which limits portability to settings where
only passive videos or observation streams are available.

Matched-baseline language refers to trainable parameter count rather
than the number of coordinates in the latent state. A $16\times16$
density matrix carries more real degrees of freedom than a length-$16$
vector hidden state, even when the trainable encoder and predictor
budgets are matched, which makes the context-probe tie in
\cref{fig:heldout} a central control: it shows that the main
comparison is not explained by a better frozen context representation.
The supplementary baseline ablations further separate target-encoder
decodability from blind-rollout performance.

The present experiments use structured partially observable oscillator
and CartPole-style settings, and the present results do not extend to
hidden Markov chains, ring-walker environments, or unstructured visual
prediction. The current comparison set covers LSTM-JEPA models with
plain and residual MLP predictors, GRU, orthogonal recurrent, and MLP
baselines, with the recurrent families grounded in
GRU~\cite{cho2014learning}, orthogonal/unitary
RNN~\cite{arjovsky2016unitary,wisdom2016full}, and
LSTM~\cite{hochreiter1997long} architectures. Learned Lie-algebra
predictors~\cite{greydanus2019hamiltonian}, neural ODE
predictors~\cite{chen2018neural}, state-space
models~\cite{gu2022efficiently,gu2023mamba}, and Transformer
predictors~\cite{vaswani2017attention} remain natural future
baselines. Within that scope, the contribution is a controlled
demonstration that belief-state structure can separate context
representation from rollout-time memory in a JEPA world model.

Seed counts are asymmetric by design. The frozen context probe is
reported over five seeds because it is the key negative control for
encoder quality; blind-rollout and action-conditioned experiments are
reported over three seeds because they require trained predictor or
counterfactual-action checkpoints. Strong claims are confined to
low-variance regions and seed counts appear in every relevant table
caption.
